% GENERATED FILE. Edit canonical lessons, then run npm run generate:books.
% canonical-source-hash: fnv1a64:90a93a6bed0f83ca
% canonical-lessons: HI-C226-sukhna, HI-C226-bigarna, HI-C226-banana, HI-C226-banna, HI-C226-tutna

\chapter{To dry, To go wrong, To make, To be made, To break}
\label{ch:hi-to-dry-to-go-wrong-to-make-to-be-made-to-break}
\hlchaptermodality{226}

\begin{chapteropening}
\textbf{By the end of this chapter:} \emph{I can say to dry, to go wrong, to make, to be made and to break.}

Complete the last lesson of chapter 226: I can say to dry, to go wrong, to make, to be made and to break.
\end{chapteropening}

\section[\texorpdfstring{\hi{सूखना} (sūkhnā)}{सूखना (sūkhnā)}]{\hi{सूखना} (sūkhnā) — to dry (by itself)}

\label{lesson:HI-C226-sukhna}



\begin{quote}
Before the new one: say the Hindi for to wring, then the Hindi for to dry.
\end{quote}

\subsection*{You'll want to know: \hi{सूखना}}
\textbf{\hi{सूखना}} — \emph{sūkhnā} — \textquotedblleft{}to dry (by itself)\textquotedblright{}.

\begin{grammarlens}[title={What you've built}]
One more word for this chapter.
\end{grammarlens}

\subsection*{Your turn}
\begin{itemize}
\raggedright
  \item \emph{Say it:} \emph{sūkhnā}
  \item \emph{Say it:} \emph{sūkhnā}, once more
  \item \emph{Say it:} say \emph{sūkhnā}
  \item \emph{Recall:} say the Hindi for to wring, then the Hindi for to dry, then say \emph{sūkhnā} again
\end{itemize}

\begin{tcolorbox}[breakable,colback=teal!4,colframe=teal!35!black,title={Before you move on}]
What does \hi{सूखना} mean? (\textquotedblleft{}To dry (by itself)\textquotedblright{}.) Say it once more.
\end{tcolorbox}
\section[\texorpdfstring{\hi{बिगड़ना} (bigaṛnā)}{बिगड़ना (bigaṛnā)}]{\hi{बिगड़ना} (bigaṛnā) — to go wrong, to spoil}

\label{lesson:HI-C226-bigarna}



\begin{quote}
Before the new one: say the Hindi for to dry, then the Hindi for to dry.
\end{quote}

\subsection*{You'll want to know: \hi{बिगड़ना}}
\textbf{\hi{बिगड़ना}} — \emph{bigaṛnā} — \textquotedblleft{}to go wrong, to spoil\textquotedblright{}.

\begin{grammarlens}[title={What you've built}]
One more word for this chapter.
\end{grammarlens}

\subsection*{Your turn}
\begin{itemize}
\raggedright
  \item \emph{Say it:} \emph{bigaṛnā}
  \item \emph{Say it:} \emph{bigaṛnā}, once more
  \item \emph{Say it:} say \emph{bigaṛnā}
  \item \emph{Recall:} say the Hindi for to dry, then the Hindi for to dry, then say \emph{bigaṛnā} again
\end{itemize}

\begin{tcolorbox}[breakable,colback=teal!4,colframe=teal!35!black,title={Before you move on}]
What does \hi{बिगड़ना} mean? (\textquotedblleft{}To go wrong, to spoil\textquotedblright{}.) Say it once more.
\end{tcolorbox}
\section[\texorpdfstring{\hi{बनाना} (banānā)}{बनाना (banānā)}]{\hi{बनाना} (banānā) — to make}

\label{lesson:HI-C226-banana}



\begin{quote}
Before the new one: say the Hindi for to dry, then the Hindi for to go wrong.
\end{quote}

\subsection*{You'll want to know: \hi{बनाना}}
\textbf{\hi{बनाना}} — \emph{banānā} — \textquotedblleft{}to make\textquotedblright{}.

\begin{grammarlens}[title={What you've built}]
One more word for this chapter.
\end{grammarlens}

\subsection*{Your turn}
\begin{itemize}
\raggedright
  \item \emph{Say it:} \emph{banānā}
  \item \emph{Say it:} \emph{banānā}, once more
  \item \emph{Say it:} say \emph{banānā}
  \item \emph{Recall:} say the Hindi for to dry, then the Hindi for to go wrong, then say \emph{banānā} again
\end{itemize}

\begin{tcolorbox}[breakable,colback=teal!4,colframe=teal!35!black,title={Before you move on}]
What does \hi{बनाना} mean? (\textquotedblleft{}To make\textquotedblright{}.) Say it once more.
\end{tcolorbox}
\section[\texorpdfstring{\hi{बनना} (bannā)}{बनना (bannā)}]{\hi{बनना} (bannā) — to be made, to become}

\label{lesson:HI-C226-banna}



\begin{quote}
Before the new one: say the Hindi for to go wrong, then the Hindi for to make.
\end{quote}

\subsection*{You'll want to know: \hi{बनना}}
\textbf{\hi{बनना}} — \emph{bannā} — \textquotedblleft{}to be made, to become\textquotedblright{}.

\begin{grammarlens}[title={What you've built}]
One more word for this chapter.
\end{grammarlens}

\subsection*{Your turn}
\begin{itemize}
\raggedright
  \item \emph{Say it:} \emph{bannā}
  \item \emph{Say it:} \emph{bannā}, once more
  \item \emph{Say it:} say \emph{bannā}
  \item \emph{Recall:} say the Hindi for to go wrong, then the Hindi for to make, then say \emph{bannā} again
\end{itemize}

\begin{tcolorbox}[breakable,colback=teal!4,colframe=teal!35!black,title={Before you move on}]
What does \hi{बनना} mean? (\textquotedblleft{}To be made, to become\textquotedblright{}.) Say it once more.
\end{tcolorbox}
\section[\texorpdfstring{\hi{टूटना} (ṭūṭnā)}{टूटना (ṭūṭnā)}]{\hi{टूटना} (ṭūṭnā) — to break (by itself)}

\label{lesson:HI-C226-tutna}



\begin{quote}
Before the new one: say the Hindi for to make, then the Hindi for to be made.
\end{quote}

\subsection*{You'll want to know: \hi{टूटना}}
\textbf{\hi{टूटना}} — \emph{ṭūṭnā} — \textquotedblleft{}to break (by itself)\textquotedblright{}.

\begin{grammarlens}[title={What you've built}]
One more word for this chapter.
\end{grammarlens}

\subsection*{Your turn}
\begin{itemize}
\raggedright
  \item \emph{Say it:} \emph{ṭūṭnā}
  \item \emph{Say it:} \emph{ṭūṭnā}, once more
  \item \emph{Say it:} say \emph{ṭūṭnā}
  \item \emph{Recall:} say the Hindi for to make, then the Hindi for to be made, then say \emph{ṭūṭnā} again
\end{itemize}

\begin{tcolorbox}[breakable,colback=teal!4,colframe=teal!35!black,title={Before you move on}]
What does \hi{टूटना} mean? (\textquotedblleft{}To break (by itself)\textquotedblright{}.) Say it once more.
\end{tcolorbox}
